\documentclass[a4paper,12pt,oneside]{article}
\makeatletter
\def\input@path{{../../}}
\makeatother
\usepackage{discrete-practice}

\begin{document}

\PracticeHeader{1}{3}{Бинарные отношения. Начало}

\Task{Задание 1}
Даны следующие бинарные отношения:
\begin{numlist}
    \item $R=\{(-1,1),(2,4),(1,1),(3,9)\}$
    \item $R=\{(a,a),(b,b),(c,c)\}$
    \item $R=\bigl\{(\{1,2\},\{\varnothing\}),(a,1),(\{1,5,a\},\{2,5,b\})\bigr\}$
\end{numlist}
Определить их всеми способами, которые вы знаете.

\Task{Задание 2}
Дополнительно указать области определения и значения.

Пусть $A=\{a,b,c,d,e\}$, а $S,T,U,V$ – отношения на $A$, где
\[
S=\{(a,a),(a,b),(b,c),(b,d),(c,e),(e,d),(c,a)\},
\]
\[
T=\{(a,b),(b,a),(b,c),(b,d),(e,e),(d,e),(c,b)\},
\]
\[
U=\{(a,b),(a,a),(b,c),(b,b),(e,e),(b,a),(c,b),(c,c),(d,d),(a,c),(c,a)\},
\]
\[
V=\{(a,b),(b,c),(b,b),(e,e),(b,a),(c,b),(d,d),(a,c),(c,a)\}.
\]

\Task{Задание 3}
Проверить тождества:
\begin{numlist}
    \item Существуют отношения, одновременно являющиеся асимметричными и несимметричными?
    \item Существуют отношения, не являющиеся симметричными и не являющиеся асимметричными?
    \item Если отношение асимметрично, то оно не является несимметричным?
    \item Если отношение не является симметричным, то оно может быть асимметричным либо несимметричным?
    \item Если отношение $aRb$ симметрично, то оно останется симметричным при перестановке элементов $a$ и $b$?
    \item Если отношение несимметрично, то оно не может быть асимметричным?
    \item Если отношение несимметрично, то оно одновременно является асимметричным?
\end{numlist}

\Task{Задание 4}
Пусть
\[
A=\{(b,a),(c,e),(d,i),(f,o),(g,u)\},\qquad
B=\{(v,a),(w,e),(x,i),(y,o),(z,u)\}.
\]
Постройте:
\begin{numlist}
    \item $A^{-1}$
    \item $B^{-1}$
    \item $A\circ A^{-1}$
    \item $A^{-1}\circ B$
    \item $A\circ B^{-1}$
\end{numlist}

\end{document}
